% Missa VIII "De Angelis" — modern notation, Latin text with Polish translation.
% Build:  lualatex -shell-escape missa-viii.tex   (needs LilyPond >= 2.24 in PATH)
\documentclass[a4paper,11pt]{article}

\usepackage[a4paper,top=15mm,bottom=18mm,left=20mm,right=20mm]{geometry}
\usepackage{fontspec}
\setmainfont{Libertinus Serif}
\usepackage[latin,main=polish]{babel}
\usepackage{microtype}
\usepackage{xcolor}
\usepackage{paracol}
\usepackage{titlesec}
\usepackage[
  program=lilypond,
  staffsize=16,
  includepaths={./,scores/},
  pass-fonts,
  tmpdir=build-ly,
]{lyluatex}

\definecolor{rubric}{RGB}{170,20,20}
\setlength{\parindent}{0pt}
\setlength{\parskip}{0pt}
\pagestyle{plain}

% \chant{Latin title}{mode}
\titleformat{\section}{\centering\Large\scshape}{}{0pt}{}
\titlespacing*{\section}{0pt}{1.8ex plus 1ex}{0.3ex}
\newcommand{\chant}[2]{%
  \section{\foreignlanguage{latin}{#1}}%
  {\centering\small\itshape #2\par}%
  \vspace{0.8ex}}

% Parallel text: Latin left, Polish right, one row per sung phrase.
\newenvironment{tekst}
  {\par\vspace{0.6ex}\small
   \setcolumnwidth{0.48\textwidth,0.48\textwidth}%
   \setlength{\columnsep}{0.04\textwidth}%
   \begin{paracol}{2}}
  {\end{paracol}\par\vspace{0.5ex}}
% \w{Latin}{Polski}
\newcommand{\w}[2]{%
  \switchcolumn[0]*\foreignlanguage{latin}{#1}\par
  \switchcolumn #2\par\vspace{0.3ex}}
% Libertinus has no U+211F/U+2123: build the barred R and V.
\newcommand{\barred}[1]{#1\kern-.48em\raisebox{.04em}{/}\kern.06em}
\newcommand{\R}{\textcolor{rubric}{\barred{R}.}~}
\newcommand{\V}{\textcolor{rubric}{\barred{V}.}~}
\newcommand{\rubr}[1]{\textcolor{rubric}{\itshape #1}}

\begin{document}

\begin{center}
  {\Huge\scshape Missa VIII}\par\smallskip
  {\Large\itshape De Angelis}\par\smallskip
  {\small Kyriale Romanum · zapis współczesny wg Vatican II Hymnal (2011)\\
   tłumaczenie polskie wg Mszału rzymskiego dla diecezji polskich}
\end{center}

% ---------------------------------------------------------------------------
\chant{Asperges me}{tonus VII}
\lilypondfile{scores/asperges.ly}
\begin{tekst}
\w{Aspérges me, Dómine, hyssópo, et mundábor:}
  {Pokrop mnie Panie hyzopem a stanę się czysty,}
\w{lavábis me, et super nivem dealbábor.}
  {obmyj mnie, a nad śnieg wybieleję.}
\w{\rubr{Ps. 50.} Miserére mei, Deus, secúndum magnam misericórdiam tuam.}
  {\rubr{Ps. 51.} Zmiłuj się nade mną, Boże, według wielkiego miłosierdzia Twego.}
\w{Glória Patri, et Fílio, et Spirítui Sancto:}
  {Chwała Ojcu, i Synowi, i Duchowi Świętemu,}
\w{Sicut erat in princípio, et nunc, et semper, et in saécula saeculórum. Amen.}
  {jak była na początku, teraz i zawsze i na wieki wieków. Amen.}
\w{\rubr{Powtarza się antyfonę} Aspérges me.}
  {\rubr{Powtarza się antyfonę} Pokrop mnie.}
\end{tekst}

% ---------------------------------------------------------------------------
\chant{Kyrie}{tonus V}
\lilypondfile{scores/kyrie.ly}
\par\smallskip
{\small\itshape\textcolor{rubric}{Jeśli Kyrie śpiewa się w tradycyjnej formie
dziewięciokrotnej, jako dziewiąte (ostatnie) wezwanie śpiewa się:}\par}
\lilypondfile{scores/kyrie-ix.ly}
\begin{tekst}
\w{Kýrie, eléison. Christe, eléison. Kýrie, eléison.}
  {Panie, zmiłuj się nad nami. Chryste, zmiłuj się nad nami.
   Panie, zmiłuj się nad nami.}
\end{tekst}

\clearpage
% ---------------------------------------------------------------------------
\chant{Gloria}{tonus V}
{\centering\small\rubr{Jeśli to możliwe, intonuje sam celebrans.}\par}
\lilypondfile{scores/gloria.ly}
\begin{tekst}
\w{Glória in excélsis Deo.}{Chwała na wysokości Bogu,}
\w{Et in terra pax homínibus bonae voluntátis.}{a na ziemi pokój ludziom dobrej woli.}
\w{Laudámus te. Benedícimus te. Adorámus te. Glorificámus te.}
  {Chwalimy Cię. Błogosławimy Cię. Wielbimy Cię. Wysławiamy Cię.}
\w{Grátias ágimus tibi propter magnam glóriam tuam.}
  {Dzięki Ci składamy, bo wielka jest chwała Twoja.}
\w{Dómine Deus, Rex caeléstis, Deus Pater omnípotens.}
  {Panie Boże, Królu nieba, Boże, Ojcze wszechmogący.}
\w{Dómine Fili unigénite, Jesu Christe.}
  {Panie, Synu Jednorodzony, Jezu Chryste.}
\w{Dómine Deus, Agnus Dei, Fílius Patris.}
  {Panie Boże, Baranku Boży, Synu Ojca.}
\w{Qui tollis peccáta mundi, miserére nobis.}
  {Który gładzisz grzechy świata, zmiłuj się nad nami.}
\w{Qui tollis peccáta mundi, súscipe deprecatiónem nostram.}
  {Który gładzisz grzechy świata, przyjm błaganie nasze.}
\w{Qui sedes ad déxteram Patris, miserére nobis.}
  {Który siedzisz po prawicy Ojca, zmiłuj się nad nami.}
\w{Quóniam tu solus sanctus. Tu solus Dóminus.}
  {Albowiem tylko Tyś jest święty. Tylko Tyś jest Panem.}
\w{Tu solus Altíssimus, Jesu Christe.}
  {Tylko Tyś Najwyższy, Jezu Chryste.}
\w{Cum Sancto Spíritu, in glória Dei Patris. Amen.}
  {Z Duchem Świętym w chwale Boga Ojca. Amen.}
\end{tekst}

\clearpage
% ---------------------------------------------------------------------------
\chant{Sanctus}{tonus VI}
\lilypondfile{scores/sanctus.ly}
\begin{tekst}
\w{Sanctus, Sanctus, Sanctus Dóminus Deus Sábaoth.}
  {Święty, Święty, Święty, Pan Bóg Zastępów.}
\w{Pleni sunt caeli et terra glória tua.}
  {Pełne są niebiosa i ziemia chwały Twojej.}
\w{Hosánna in excélsis.}{Hosanna na wysokości.}
\w{Benedíctus qui venit in nómine Dómini.}
  {Błogosławiony, który idzie w imię Pańskie.}
\w{Hosánna in excélsis.}{Hosanna na wysokości.}
\end{tekst}

% ---------------------------------------------------------------------------
\chant{Agnus Dei}{tonus VI}
\lilypondfile{scores/agnus.ly}
\begin{tekst}
\w{Agnus Dei, qui tollis peccáta mundi: miserére nobis.}
  {Baranku Boży, który gładzisz grzechy świata, zmiłuj się nad nami.}
\w{Agnus Dei, qui tollis peccáta mundi: miserére nobis.}
  {Baranku Boży, który gładzisz grzechy świata, zmiłuj się nad nami.}
\w{Agnus Dei, qui tollis peccáta mundi: dona nobis pacem.}
  {Baranku Boży, który gładzisz grzechy świata, obdarz nas pokojem.}
\end{tekst}

% ---------------------------------------------------------------------------
\chant{Ite, missa est}{tonus V}
\lilypondfile{scores/ite.ly}
\begin{tekst}
\w{\V Ite, missa est.}{\V Idźcie w pokoju Chrystusa.}
\w{\R Deo grátias.}{\R Bogu niech będą dzięki.}
\end{tekst}

\vfill
{\centering\itshape + Ad maiorem Dei gloriam +\par}

\end{document}
